\documentclass[11pt]{amsart}
\usepackage[margin=1in]{geometry}
\usepackage{amssymb,amsmath,amsthm,mathtools}
\usepackage{booktabs,array}
\usepackage[hidelinks]{hyperref}
\hypersetup{pdftitle={Forced-Entry Costs at a Fixed Optimal Assignment under Paired Interval Uncertainty},pdfauthor={Filip Majchrzak}}
\IfFileExists{microtype.sty}{\usepackage[expansion=false]{microtype}}{}
\emergencystretch=3em

\theoremstyle{plain}
\newtheorem{theorem}{Theorem}
\newtheorem{proposition}[theorem]{Proposition}
\newtheorem{lemma}[theorem]{Lemma}
\newtheorem{corollary}[theorem]{Corollary}
\theoremstyle{definition}
\newtheorem{definition}[theorem]{Definition}
\theoremstyle{remark}
\newtheorem{remark}[theorem]{Remark}

\newcommand{\R}{\mathbb{R}}
\newcommand{\Z}{\mathbb{Z}}
\newcommand{\E}{\mathbb{E}}
\newcommand{\one}{\mathbf{1}}
\newcommand{\cA}{\mathcal{A}}
\newcommand{\COR}{\mathrm{COR}}
\newcommand{\xc}{\operatorname{xc}}
\newcommand{\size}{\operatorname{size}}
\newcommand{\rk}{\operatorname{rank}_+}
\newcommand{\Ent}{\mathsf{H}}
\newcommand{\MI}{\mathsf{I}}
\newcommand{\TV}{\operatorname{TV}}
\newcommand{\code}[1]{\texttt{#1}}

\title[Forced-entry costs under paired interval uncertainty]{Forced-Entry Costs at a Fixed Optimal Assignment\\ under Paired Interval Uncertainty:\\ a Lean-Verified Exact Reduction\\ and Its Consequences for Linear Programs}
\author{Filip Majchrzak}
\address{Independent researcher, Poland (SAPIO MORPHEUS project)}
\thanks{ORCID 0009-0007-3475-8900. Version 1.0, 24 November 2026. Archived at Zenodo, \href{https://doi.org/10.5281/zenodo.24511022}{doi:10.5281/zenodo.24511022}. Text licensed CC BY 4.0; code Apache-2.0.}
\date{24 November 2026}

\begin{document}

\begin{abstract}
Let the cost matrix of an assignment problem depend on interval parameters, and suppose that one assignment $F$ is
optimal in every scenario. For an entry $e\notin F$, the smallest extra cost, over all scenarios, of forcing $e$
into the solution is a tolerance-type quantity of sensitivity analysis. When the costs are nondecreasing in the parameters and no parameter is shared between
an entry of $F$ and an entry outside $F$, it is obtained from a single assignment problem. We show that the simplest sharing changes this completely,
even in a very restricted model: $F$ is the unique optimum in every scenario and has an explicit dual certificate,
every parameter occurs in exactly two entries with coefficient $+1$, the uncertain entries form a matching, and all
costs are nonnegative. An explicit construction, without perturbation, turns every quadratic shortest-path instance on a directed
acyclic graph with nonnegative coefficients into such an instance: in every scenario the forcing cost is, up to a
fixed additive constant, the value of a linear shortest-path problem on the same graph, and its minimum over the
scenarios is the quadratic optimum. Hence the quantity is NP-hard, and for general instances of the model it cannot
be approximated within a factor $2^{\mathrm{poly}(N)}$ unless $\mathrm{P}=\mathrm{NP}$. We then record what classical results on extension complexity
imply for linear programs with a fixed feasible region and an objective affine in the weights. For an explicit
family of $N\times N$ instances, $N=8n^2-n+1$, such a program needs exactly as many inequalities as an extended
formulation of the correlation polytope $\COR(n)$, hence at least $1.5^n$ (Kaibel and Weltge), also for weights in
$\{1,\dots,5\}$; for interactions along the edges of a graph $G$ the exact size is $\xc(\COR(G))$, which is
$2^{O(\operatorname{tw}(G)+\log n)}$ with a matching lower bound in proper minor-closed classes (Aboulker et al.) and
is $2^{\Omega(n/\log n)}$, in dimension $O(n)$, for an explicit family of bounded-degree graphs (G\"o\"os, Jain and
Watson); every fixed additive error still needs $2^{\Omega(n)}$ inequalities, while for weights in $[1,5]$ every fixed
relative error is achieved with $n^{O(\varepsilon^{-2})}$ inequalities. The largest forcing cost over the scenarios,
by contrast, has a representation with $4n^2$ inequalities. The exponential bounds and the rounding
argument are known. What is added is the reduction inside the restricted model and elementary lemmas that recover
lifts of solutions from optimal values alone.
The reduction theorem is formalised in full in Lean~4 with Mathlib; the statements on linear programs and the corollaries are not machine-checked, apart from the covering inequality of Kaibel and Weltge.
\end{abstract}

\maketitle

\section{Introduction}

\subsection*{The question}
Let $C(\theta)$ be an $N\times N$ cost matrix depending on a vector of parameters $\theta$ that ranges over a box
$\Theta$, and let $\cA_N$ be the polytope of doubly stochastic matrices. Suppose that one permutation matrix $F$ is
optimal for $\min\{C(\theta)\cdot X:X\in\cA_N\}$ in \emph{every} scenario $\theta\in\Theta$. For an entry $e$ with
$F_e=0$, the \emph{forcing cost}
\[
 h_e(\theta):=\min\{\,C(\theta)\cdot(X-F)\;:\;X\in\cA_N,\ X_e=1\,\}\ \ge0
\]
measures how much is lost in scenario $\theta$ if $e$ must be used, and its \emph{lower endpoint}
\[
 \eta_e:=\min_{\theta\in\Theta}h_e(\theta)
\]
is the amount by which the cost of $e$ has to decrease before $e$ can belong to an optimal assignment in some
scenario. For a single cost matrix this is the classical lower tolerance of an element outside the optimal solution
\cite{GJM,LinWen}; with interval data it is a quantity of the kind studied in tolerance analysis and in interval
combinatorial optimisation \cite{Wendell,Hladik11,Hladik26,KZ}. In the vocabulary of interval optimisation, the
hypothesis says that $F$ is \emph{necessarily optimal} (see \cite{Hladik12} for the complexity of this property in
interval linear programming), and $\eta_e>0$ says that $e$ is not \emph{possibly optimal}.

\begin{proposition}[no parameter shared across $F$]\label{prop:indep}
Let $\Theta$ be a box, let $C(\theta)$ be affine in $\theta$ with nonnegative coefficients, and suppose that every
parameter occurs only in entries of $F$ or only in entries outside $F$. Then $\eta_e=h_e(\theta^\ast)$, where
$\theta^\ast$ puts the parameters occurring in $F$ at their upper and all other parameters at their lower endpoints.
So $\eta_e$ is obtained by solving one assignment problem. In particular this holds when every entry of $C$ varies
independently in an interval.
\end{proposition}
\begin{proof}
By Birkhoff's theorem the minimum defining $h_e(\theta)$ is attained at a permutation matrix $X$, and
$C(\theta)\cdot(X-F)=\sum_{X\setminus F}C_{ij}(\theta)-\sum_{F\setminus X}C_{ij}(\theta)$ is nondecreasing in the
parameters that occur outside $F$ and nonincreasing in those that occur in $F$. So it is smallest at
$\theta^\ast$, for all $X$ simultaneously.
\end{proof}
The same argument shows that $F$ is optimal in every scenario if and only if it is optimal at $\theta^\ast$. The
proposition is elementary and presumably folklore (extreme scenarios of this kind are standard in optimisation with
interval costs); we did not find a source.

Thus, when all coefficients are nonnegative, a parameter can cause difficulty only if it is shared by an entry of
$F$ and an entry outside $F$. The sign condition is needed: coefficients $+1$ and $-1$ in two entries outside $F$
produce the same product term as in Theorem~\ref{thm:adapter} below (put $\beta_j-\theta_j$ on the arc $g_j\to g_j'$
of Section~\ref{sec:adapter} and $0$ on the diagonal). We ask what happens in the mildest case with nonnegative
coefficients: each parameter occurs in exactly \emph{two} entries, one on each side.

\subsection*{The restricted model}
We impose the following conditions, which make the instance as benign as we know how to make it.
\begin{itemize}
\item[(C1)] \emph{Fixed certified optimum.} $F$ is the identity and is the unique optimal point of $\cA_N$ for every
  $\theta\in\Theta$; an optimal dual solution is given explicitly and is affine in $\theta$.
\item[(C2)] \emph{Paired uncertainty.} $\Theta=\prod_j[0,\beta_j]$, and each $\theta_j$ occurs in exactly two entries
  of $C(\theta)$, with coefficient $+1$ in both; one of the two entries lies on $F$.
\item[(C3)] \emph{Matching support.} No two uncertain entries share a row or a column.
\item[(C4)] \emph{Nonnegativity.} All entries of $C(\theta)$ are nonnegative for all $\theta\in\Theta$.
\end{itemize}
Throughout, a \emph{scenario} is a point of the box $\Theta$; the box is the whole uncertainty set.

\subsection*{The reduction}
Let $D=(V,A)$ be a directed acyclic graph with terminals $s\ne t$, arc costs $\alpha_a\ge0$, and a list $J$ of
\emph{interactions}: ordered pairs $(a_j,b_j)$ of distinct arcs with coefficients $\beta_j\ge0$. The cost of an
$s$--$t$ path $P$ is $Q(P)=\sum_{a\in P}\alpha_a+\sum_{j:\,a_j,b_j\in P}\beta_j$; minimising $Q$ is the quadratic
shortest path problem \cite{Rostami,HuSotirov}.

\begin{theorem}[exact adapter]\label{thm:adapter}
For every such instance with at least one $s$--$t$ path and every $\varrho>0$ there is an explicit complete
assignment instance of dimension
$N=|V|+|A|+4|J|$, with parameters $\theta_j\in[0,\beta_j]$ $(j\in J)$ and observed entry $e=(t,s)$ of cost $\varrho$,
which satisfies \textup{(C1)--(C4)} and
\[
 h_e(\theta)=\varrho+\min_{P}\Bigl[\sum_{a\in P}\alpha_a
      +\sum_{j\in J}\bigl(\theta_j\,[a_j\in P]+(\beta_j-\theta_j)\,[b_j\in P]\bigr)\Bigr]
 \qquad\text{for every }\theta\in\Theta ,
\]
the minimum being over the $s$--$t$ paths of $D$. Consequently $\eta_e=\varrho+\min_PQ(P)$.
\end{theorem}

The last sentence rests on the elementary identity $\min_{0\le\theta\le\beta}\{\theta u+(\beta-\theta)v\}=\beta uv$ for
$u,v\in\{0,1\}$. If $\beta_j=0$ the interval of $\theta_j$ is a point and \textup{(C2)} holds only formally; this
occurs in Theorem~\ref{thm:exact} and Corollary~\ref{cor:tw} below, where zero weights are allowed, but not in Corollaries \ref{cor:np} and
\ref{cor:inapx} or in Theorems \ref{thm:bounded}--\ref{thm:rel}. The construction uses no perturbation and no
limiting argument. Its ingredients are classical: a path of a directed acyclic graph as the unique cycle through a
return arc in an assignment (the link between shortest paths and assignments usually attributed to \cite{HM}; see
\cite[Section~IV-A]{MS} for a recent use on general graphs), and the linearisation of a product of binary variables
\cite{GW}. What we have not found in the literature is the
combination that realises the pairwise terms by \emph{one interval parameter shared by an entry of $F$ and an entry
outside $F$} while keeping $F$ optimal, with a certificate, in every scenario.

For $n\ge1$ let $D_n$ be a chain of $n$ ``diamonds'': an $s$--$t$ path chooses one of two branches in each layer
and so corresponds to $z\in\{0,1\}^n$. Put a weight $w_{i,a}$ on branch $a$ of layer $i$ and an interaction of
weight $w_{ij,ab}$ between branch $a$ of layer $i$ and branch $b$ of layer $j$, for all $i<j$ and $a,b\in\{0,1\}$.
The weight vector $w$ has $d_n=2n^2$ coordinates, and with $\varrho=1$ Theorem~\ref{thm:adapter} gives a family of
assignment instances of dimension $N=8n^2-n+1$ whose combinatorial data (dimension, $F$, $e$, positions of the
parameters) do not depend on $w$, with
\begin{equation}\label{eq:eta}
 \eta_n(w)=1+\min_{z\in\{0,1\}^n}w\cdot m_n(z)\qquad(w\in\R^{d_n}_{\ge0}),
\end{equation}
where $m_n(z)\in\{0,1\}^{d_n}$ lists the indicators $[z_i=a]$ and $[(z_i,z_j)=(a,b)]$.

\begin{corollary}\label{cor:np}
Computing $\eta_e$ is NP-hard for instances satisfying \textup{(C1)--(C4)}, already for the family
\eqref{eq:eta} with weights in $\{1,5\}$.
\end{corollary}

Together with Proposition~\ref{prop:indep} this is a sharp contrast: parameters with nonnegative coefficients that
stay on one side of $F$ are harmless however often they occur, and one parameter pattern, an entry of $F$ sharing its parameter with one entry
outside $F$, already makes the quantity NP-hard, in a model where recognising the optimum is not the difficulty.

\begin{corollary}[no approximation in general]\label{cor:inapx}
Let $p$ be a polynomial with nonnegative integer coefficients. Unless $\mathrm{P}=\mathrm{NP}$, there is no
polynomial-time algorithm that, for every
instance satisfying \textup{(C1)--(C4)} with integer data and dimension $N$, returns a number between $\eta_e$ and
$2^{p(N)}\,\eta_e$.
\end{corollary}

This transports to the restricted assignment model the known fact that the quadratic shortest path problem cannot
be approximated \cite{Rostami}. It concerns general directed acyclic graphs with few interactions and is proved
from the problem of paths avoiding forbidden pairs \cite{GMO}; for the dense positive family \eqref{eq:eta} the
picture is different (Theorem~\ref{thm:rel} below).

\subsection*{Consequences for linear programs}
The remaining results are unconditional statements about one model of linear programs. They follow from classical
theorems on extension complexity once a few short lemmas are in place; we state plainly which part is which.

\begin{definition}\label{def:lp}
Let $W\subseteq\R^d_{\ge0}$ and $\eta:W\to\R$. A \emph{weight-affine LP representation} of $\eta$ on $W$ consists of a
nonempty polyhedron $P\subseteq\R^k$, an affine function $a:\R^k\to\R$ and an affine map $\pi:\R^k\to\R^d$, all
independent of $w$, with
\[
 \min_{y\in P}\{a(y)+w\cdot\pi(y)\}=\eta(w)\qquad\text{for all }w\in W .
\]
We call $(P,a,\pi)$ a \emph{weight-affine LP} and $w\mapsto\min_{y\in P}\{a(y)+w\cdot\pi(y)\}$ its \emph{value
function}. Its \emph{size}, $\size(P)$, is the smallest number of inequalities in a description of $P$ by linear
equations and inequalities; equations are free. $P$ may be unbounded and of any dimension.
\end{definition}

Thus the feasible region is fixed and the objective is affine in the weights themselves. This is close to the
non-uniform model of formulation complexity of Braun, Pokutta and Zink \cite{BPZ}, with two differences that point
in opposite directions. We require the objective to depend on the weights through one fixed affine pair $(a,\pi)$,
whereas in \cite{BPZ} every instance has its own affine objective. And we do not assume that feasible solutions are
realised as points of $P$ on which the objectives take their values: only the optimal values are prescribed.
For exact representations of $\eta_n$ on the weight sets of Theorems \ref{thm:exact} and \ref{thm:bounded},
Lemma~\ref{lem:face} and Step~1 of the proof of Theorem~\ref{thm:bounded} supply the realisations, after which such
a representation is a formulation in the sense of \cite{BPZ}; in Theorem~\ref{thm:abs} only the approximate lifts
\eqref{eq:lift} are available. Definitions through values appear in
\cite[Section~2.2]{BFPS}, where the lower inequality is required for all objectives, which amounts to containing
the polytope, and in \cite{HL}. For a polytope $Q$, $\xc(Q)$ denotes the smallest number of inequalities of a polyhedron having $Q$ as
an affine image.

\begin{theorem}[exact size]\label{thm:exact}
The smallest size of a weight-affine LP representation of $\eta_n$ on $W=\Z^{d_n}_{\ge0}$ is exactly
$\xc(\COR(n))$. In particular it lies between $1.5^n$ and $2^n$.
\end{theorem}

\begin{theorem}[bounded positive weights]\label{thm:bounded}
Every weight-affine LP representation of $\eta_n$ on $W=\{1,2,3,4,5\}^{d_n}$ has size at least $1.5^n$.
\end{theorem}

In terms of the assignment dimension this is $2^{\Omega(\sqrt N)}$; for exact representations on nonnegative
integer weights and another explicit family, Corollary~\ref{cor:tw} gives $2^{\Omega(N/\log N)}$. The number $1.5^n$ is the
bound of Kaibel and Weltge \cite{KW} for $\COR(n)$ and for coverings of the unique-disjointness matrix; the
exponential lower bound for $\COR(n)$ goes back to Fiorini, Massar, Pokutta, Tiwary and de Wolf \cite{FMPTW}. Our
contribution to these two theorems is Lemma~\ref{lem:face} and Step~1 of the proof of Theorem~\ref{thm:bounded},
which recover exact lifts of the vertices from optimal values (in Step~1, on finitely many weight vectors). For weights in
$\{1,\dots,5\}$ every parameter interval has positive length, all costs are bounded by a polynomial in $n$, and the
completion cost can be taken constant (Remark~\ref{rem:const}).

The dense pattern of interactions is not essential for the exact statement. For a graph $G$ on $[n]$ keep only the
interactions between layers $i<j$ with $ij\in E(G)$. Let $\eta_G$ be the resulting function of the $2n+4|E(G)|$
weights, and let $\COR(G)=\operatorname{conv}\{((z_i)_{i\in[n]},(z_iz_j)_{ij\in E(G)}):z\in\{0,1\}^n\}$ be the
correlation polytope of $G$, so that $\COR(n)=\COR(K_n)$.

\begin{corollary}[sparse interactions]\label{cor:tw}
The instances defining $\eta_G$ have dimension $N_G=7n+1+16|E(G)|$ and satisfy \textup{(C1)--(C4)}, and the smallest
size of a weight-affine LP representation of $\eta_G$ on nonnegative integer weights is exactly $\xc(\COR(G))$.
Consequently, by the theorems of Aboulker, Fiorini, Huynh, Macchia and Seif \cite{AFHMS}, it is
$2^{O(\operatorname{tw}(G)+\log n)}$ for every $G$, where $\operatorname{tw}(G)$ is the treewidth of $G$, and at
least $2^{c(\operatorname{tw}(G)+\log n)}$ for all $G$ in a fixed proper minor-closed class, with $c>0$ depending
on the class. Moreover, for the explicit bounded-degree graphs of G\"o\"os, Jain and Watson \cite{GJW} it is
$2^{\Omega(n/\log n)}=2^{\Omega(N_G/\log N_G)}$.
\end{corollary}

So the size is polynomial when the interaction graph has bounded treewidth, and for the $h\times h$ grid, where
$n=h^2$ and $N_G=\Theta(n)$, it is $2^{\Theta(\sqrt{N_G})}$. The last statement of the corollary improves the
dependence on the assignment dimension from $2^{\Omega(\sqrt N)}$ for the dense family to $2^{\Omega(N/\log N)}$; a
bound $2^{\Omega(N_G)}$ would follow for all graphs with $O(n)$ edges and treewidth $\Omega(n)$ from
\cite[Conjecture~2]{AFHMS}, which asserts the lower bound for all graphs, and for the graphs of \cite{GJW} from
their conjecture that the logarithm in their bound can be removed.

The difficulty belongs to the lower endpoint only. As a minimum of functions affine in $\theta$, $h_e$ is concave;
minimising it over the box is the hard direction, maximising it is a linear program.

\begin{proposition}[the upper endpoint]\label{prop:upper}
Let $C(\theta)$ be affine in $\theta$ and let $\Theta$ be a polytope described by $k$ inequalities. Then
$\max_{\theta\in\Theta}h_e(\theta)$ is the optimal value of a linear program with $(N-1)^2+k$ inequalities. For the
instances of Theorem~\ref{thm:adapter},
\[
 \max_{\theta\in\Theta}h_e(\theta)=\varrho+\max\bigl\{\pi_t-\pi_s:\ \pi_y-\pi_x\le\ell_a(\theta)\ \text{for every arc }
 a=(x,y)\text{ of }D,\ 0\le\theta\le\beta\bigr\},
\]
where $\ell_a(\theta)=\alpha_a+\sum_{j:\,a_j=a}\theta_j+\sum_{j:\,b_j=a}(\beta_j-\theta_j)$ and $\pi$ ranges over the
real functions on the vertices of $D$. For the family \eqref{eq:eta}, the function
$w\mapsto\max_{\theta}h_e(\theta)$ has a weight-affine LP representation with at most $4n^2$ inequalities on
$\R^{d_n}_{\ge0}$.
\end{proposition}

So in one and the same model the two endpoints of the interval of forcing costs behave in opposite ways: $4n^2$
inequalities suffice for the upper endpoint, while the lower endpoint needs at least $1.5^n$. The proposition is an
application of linear programming duality of a kind that is standard in robust optimisation; we state it because
it settles which endpoint carries the difficulty.

\begin{theorem}[fixed additive error]\label{thm:abs}
There are absolute constants $c>0$ and $n_0$ such that the following holds. Let $\Delta\ge0$, and let $g$ be the
value function of a weight-affine LP $(P,a,\pi)$ with $|g(w)-\eta_n(w)|\le\Delta$ for all
$w\in\{1,\dots,5\}^{d_n}$. Put $\nu=\max\{1,\lceil40\Delta\rceil\}$ and $m=\lfloor n/\nu\rfloor$. If $m\ge n_0$, then $P$
has at least $2^{cm}$ inequalities.
\end{theorem}

\begin{theorem}[fixed relative error]\label{thm:rel}
For integers $1\le t<n$ there is a weight-affine LP with at most $2^{t+1}\binom n{t+1}$ inequalities whose value
$L_t$ satisfies, for all real $w\in[1,5]^{d_n}$,
\[
 1+\kappa_n\le L_t(w)\le\eta_n(w)\le L_t(w)+n^2\sqrt{2\ln2/t},\qquad \kappa_n=\tbinom{n+1}{2}.
\]
Hence for $0<\varepsilon\le1$ and $t=\lceil8\ln2/\varepsilon^2\rceil$ (using the exact $2^n$-inequality
representation if $t\ge n$) one has $L_t\le\eta_n\le(1+\varepsilon)L_t$ with $\min\{2^n,\,n^{O(\varepsilon^{-2})}\}$
inequalities.
\end{theorem}

Theorem~\ref{thm:abs} applies the corruption lemma of Braun, Fiorini, Pokutta and Steurer \cite{BFPS}, whose
Theorem~6 already gives $2^{\Omega(n)}$ for approximations of $\COR(n)$ that contain it; here the lifts are only
approximate, and a further elementary estimate, \eqref{eq:lift}, controls them. Theorem~\ref{thm:rel} is a
specialisation of correlation rounding for dense instances \cite{RT,YZ,MR,JKR}. The entropy argument is that of
\cite[Lemma~4.5]{RT}, stated there for semidefinite hierarchies; we use only the pairwise case, which is
self-contained and does not depend on the general statements discussed in \cite{JKR}. Together the two theorems
say that the obstruction of Theorem~\ref{thm:bounded} concerns \emph{exact and additively accurate} weight-affine
programs only: in this dense positive family a fixed percentage error is cheap, because
$\eta_n(w)\ge1+\kappa_n=\Theta(n^2)$. A percentage error that tends to zero is not cheap: if the value function $L$
of a weight-affine LP satisfies $L\le\eta_n\le(1+\varepsilon)L$ on $\{1,\dots,5\}^{d_n}$ with
$\varepsilon\le1/(40(1+5\kappa_n))$, then $0\le\eta_n-L\le\varepsilon\eta_n\le\frac1{40}$, and
Theorem~\ref{thm:abs} with $\nu=1$ gives at least $2^{cn}$ inequalities for $n\ge n_0$, while for
$\varepsilon=1/(40(1+5\kappa_n))$ one has $1/\varepsilon=100n(n+1)+40$. So no family of such programs has size polynomial in $n$ and $1/\varepsilon$; compare the
polynomial-size and fully polynomial-size relaxation schemes of \cite{BPmatching}. The constant $c$ in
Theorem~\ref{thm:abs} is small (at most $1/(256\ln2)$), the two theorems address two extreme regimes, and the
range between them is open.

\subsection*{What is and is not shown}
The lower bounds concern one fixed polyhedron with an objective affine in $w$ itself. They do not bound the running
time of algorithms and do not exclude: a different linear program for each $w$; programs in which $w$ enters the
constraint matrix; semidefinite programs; nonlinear preprocessing of $w$. Programs in which $w$ enters only the
right-hand side are covered by duality (Remark~\ref{rem:rhs}). Two caveats deserve emphasis. First, the model is
not invariant under re-encoding of the weights. Avis, Bremner, Tiwary and Watanabe \cite{ABTW} show that every
decision problem in P/poly has polynomial-size \emph{weak} extended formulations, in which the objective is affine
in a $0/1$ encoding of the input; for a two-letter alphabet the weight vector is such an encoding. We therefore
expect (the inference is ours, and it needs their construction for the $O(\log n)$ output bits of $\eta_n$, which we
have not written out) that a super-polynomial bound for two letters would show that $\eta_n$ has no polynomial-size
circuits there; as $\eta_n$ can be computed in polynomial time with an NP oracle, this would show
$\mathrm{NP}\not\subseteq\mathrm{P/poly}$. So the NP-hardness of Corollary~\ref{cor:np} for the alphabet $\{1,5\}$ comes with no unconditional bound
on linear programs; Theorem~\ref{thm:bounded} is proved with four letters (Remark~\ref{rem:letters}), and three letters are
open. Second, representing a function as a difference of two linear programs can be exponentially cheaper
\cite{BHKL}. The
theorems say that an easy feasible region, a fixed optimal solution, an explicit dual certificate and a sparse
pattern of dependence do not by themselves give a compact linear description of the lower endpoint $\eta_e$ as a
function of the data, although they do give one of the upper endpoint (Proposition~\ref{prop:upper}).

\subsection*{Related work}
\emph{Sensitivity and tolerance analysis.} The tolerance approach goes back to Wendell \cite{Wendell}. For a fixed
cost vector, tolerances and stability radii are classical notions \cite{SLG,LinWen,GJM}, and for polynomially
solvable problems they can be computed in polynomial time \cite{CW}. For general linear programs, determining the maximal tolerances of a possibly degenerate
optimal solution is NP-hard \cite{Hladik11}; Hlad\'{\i}k \cite{Hladik26} strengthens this to inapproximability,
gives bounds, and extends the approach to linear dependencies among the cost coefficients. Our model is a special
case of such dependencies, for the assignment problem, with a fixed box and a different question (one direction
instead of a common scale). We could read only the abstract and the reference list of \cite{Hladik26}; neither
indicates an assignment-specific reduction, a two-occurrence restriction or statements on LP size, but we cannot
exclude an overlap. For interval weights, deviation intervals of elements and their complexity are studied in
\cite{KZ}, and the corresponding questions for activity networks in \cite{Fortin}; we know both from summaries only
and cannot say how close their hardness results come to Corollary~\ref{cor:np}. Min-max and min-max regret
assignment problems \cite{ABV} ask a different question.

\emph{Quadratic shortest paths and conflicts.} The quadratic shortest path problem is NP-hard and cannot be
approximated unless $\mathrm{P}=\mathrm{NP}$ \cite{Rostami}; Hu and Sotirov reduce the quadratic assignment problem
to it on layers of parallel arcs \cite{HuSotirov}, the same kind of topology as our diamond chain, which has two
arcs per layer; see \cite{Cela} for linearisable instances. Finding a path that avoids forbidden pairs of vertices
is NP-hard on directed acyclic graphs \cite{GMO,Kovac}. Matching problems become hard under pairwise conflicts even
when every component of the conflict graph is a single edge \cite{DPSW}.

\emph{Extension complexity.} See \cite{Yannakakis,FMPTW,KW} for exact and \cite{BFPS,CLRS,BPZ} for approximate
formulations. Bounds for approximations of the correlation polytope whose size depends on the approximation
parameter are attributed in \cite[Example~3.10]{BPZ} to Braverman and Moitra \cite{BM} and to Braun and Pokutta
\cite{BP13}; the dependence $2^{\Omega(n/\Delta)}$ in Theorem~\ref{thm:abs}, obtained here by a copy argument, should
be compared with those results, which we have not read. The equality of the extension complexity of a polytope and of the epigraph of its support function,
and value-only approximations for nonnegative objectives, are in \cite{HL}. We did not find Lemma~\ref{lem:face} or
the estimate \eqref{eq:lift} stated in the form used here; both are short.

\subsection*{Notation}
$[\,\cdot\,]$ is the Iverson bracket, $C\cdot X=\sum_{u,v}C_{uv}X_{uv}$, $\one$ is the all-ones vector and $\one_T$
the indicator vector of a set $T$. Several letters are reused. In the statements on the reduction and in Sections
\ref{sec:adapter} and \ref{sec:diamond}, $P$ is a path, $a$ and $b$ are arcs or branch indices, $s$ and $t$ are the
terminals, and $e$ is the observed entry. In Definition~\ref{def:lp}, in the statements on linear programs and in
Sections \ref{sec:exact}--\ref{sec:rel}, $P$ is a polyhedron, $a$ and $b$ are affine functions (in Remark~\ref{rem:rhs}, $b$ is
a vector), and $S,T$ are sets or tuples of indices; in Section~\ref{sec:rel} the letters $r,s,t$ are integers.

\subsection*{Status of verification}
This work was produced by AI agents under the author's direction (see the end of the paper). We therefore state what
has been checked and how.
\begin{itemize}
\item \emph{Proof review.} The proofs were first written as project notes; a separate AI agent of the same project
  re-derived them step by step and reported no error. The present text was then reviewed by another AI agent, which
  re-derived Sections \ref{sec:adapter}--\ref{sec:rel} as written, implemented the construction of
  Section~\ref{sec:adapter} from this text alone, and compared the two imported statements, \cite[Lemma~4]{BFPS}
  and the covering bound of \cite{KW}, with their arXiv versions. It found no error in a theorem, proof step,
  constant or count. It did find inaccurate sentences about cited work and about what had been checked, a hypothesis
  missing from the statement of Theorem~\ref{thm:adapter} in the introduction, an open question that is settled by
  duality (now Remark~\ref{rem:rhs}), and slips in two side remarks; all are corrected here. The passages added or changed after that review, among them Corollaries \ref{cor:inapx} and \ref{cor:tw}, Remark~\ref{rem:rhs}, Proposition~\ref{prop:upper}, the general form of Proposition~\ref{prop:indep} and Section~\ref{sec:formal}, were checked in three further passes by other AI agents, two of them with their own code; they found no error and several formulations that claimed too much, corrected here.
\item \emph{Computation.} The construction was implemented independently four times: once from the project's
  written description and three times from Section~\ref{sec:adapter} of this text. The numbers below are those of
  the first two implementations. In exact rational arithmetic, for the diamond family
  with $n\le3$ and for $96$ random directed acyclic graphs, in corner scenarios and in sampled interior ones, the
  identity is optimal, the dual certificate is feasible, and the forced value equals the route minimum; uniqueness
  of the optimal permutation was checked for $n=2,3$ in a subsample of the scenarios, for the random graphs, and by
  enumeration of all permutations for instances with $N\le9$. For $n=4,5$ $(N=125,196)$ optimality, the dual
  certificate and the forced value were checked in sampled scenarios, with a floating-point assignment solver on
  integer-scaled matrices. Formula \eqref{eq:eta} was confirmed for $n=2,3$ over all corner scenarios, the formula in
  the proof of Corollary~\ref{cor:np} for all graphs on at most three vertices, and the reduction in the proof of
  Corollary~\ref{cor:inapx} on $250$ random instances (two implementations). Proposition~\ref{prop:indep} was
  checked by enumeration on $400$ random instances with $N\le4$, Corollary~\ref{cor:tw} (dimension and value formula)
  on $10$ instances with $n\le4$. With a floating-point linear programming solver: the inequalities of
  Theorem~\ref{thm:rel} for $n\le6$; the two programs in the proof of Proposition~\ref{prop:upper} on $40$ random
  graphs, with the exact forced value at the maximising scenario; and, for $n=2,3$, the fixed polyhedron with $4n^2$
  inequalities. These computations are consistency checks, not proofs.
\item \emph{Formal verification.} Theorem~\ref{thm:adapter} is machine-checked in Lean~4 in full: the dimension,
  conditions \textup{(C1)--(C4)} including uniqueness among all doubly stochastic matrices, the formula for
  $h_e(\theta)$ in every scenario and the formula for $\eta_e$; see Section~\ref{sec:formal} for the exact statement
  and its qualifications. The covering inequality of \cite{KW} is formalised as well. Not formalised:
  Section~\ref{sec:diamond}, Propositions \ref{prop:indep} and \ref{prop:upper}, the corollaries,
  Lemma~\ref{lem:face} and Theorems \ref{thm:exact}--\ref{thm:rel}.
\item \emph{Literature.} The search was limited to open sources. Of the works cited,
  \cite{AFHMS,ABTW,BHKL,BFPS,BPmatching,BPZ,Cela,CLRS,FMPTW,GJW,HL,HuSotirov,JKR,KW,Kovac,MR,MS,RT} were opened in their arXiv
  versions and the passages used were read; theorem and section numbers refer to those versions.
  \cite{ABV,CW,DPSW,Fortin,GJM,Hladik11,Hladik12,Hladik26,KZ,LinWen,Rostami,YZ} are known to us only from abstracts
  or database summaries, and \cite{BM,BP13,GJS,GMO,GW,HM,Karp,SLG,Wendell,Yannakakis} only from bibliographic
  records; statements about these two groups repeat the summaries or the standard attribution. For the main claim
  we could not read the full texts of \cite{Hladik26}, \cite{KZ} and \cite{Fortin}; see the paragraph on related
  work.
\end{itemize}
No external expert has reviewed this paper. We welcome corrections.

\section{The adapter}\label{sec:adapter}

\subsection{Construction}
Let $D,s,t,\alpha,J,\beta$ be as in the introduction, with at least one $s$--$t$ path; parallel arcs are distinct
arcs. Write $A_0=\sum_a\alpha_a$, $B_0=\sum_j\beta_j$, fix $\varrho>0$ and a number
\begin{equation}\label{eq:Gamma}
 \Gamma\ \ge\ \varrho+A_0+2B_0+1 .
\end{equation}
Replace every arc $a=(x,y)$ of $D$ by a directed chain from $x$ to $y$ through new vertices, as follows. The first
arc of the chain goes from $x$ to a new vertex $p_a$ and has cost $\alpha_a$. Then, in a fixed order, for every
interaction $j$ in which $a$ occurs:
\begin{itemize}
\item if $a=a_j$: add new vertices $f_j,f_j'$, an arc of cost $0$ from the current end of the chain to $f_j$, and an
  arc $f_j\to f_j'$ of cost $\theta_j$;
\item if $a=b_j$: add new vertices $g_j,g_j'$, an arc of cost $0$ from the current end of the chain to $g_j$, and an
  arc $g_j\to g_j'$ of cost $\beta_j$.
\end{itemize}
Finally join the current end of the chain to $y$ by an arc of cost $0$. Let $D'$ be the resulting graph; it has
$N=|V|+|A|+4|J|$ vertices. Add the \emph{return arc} $e=(t,s)$ of cost $\varrho$.

Rows and columns of the cost matrix $C(\theta)$ are indexed by the vertices of $D'$. Every arc of $D'$ and the arc
$e$ give an off-diagonal entry with the stated cost; these are the \emph{designated} entries. Diagonal entries are
$0$, except $C_{g_jg_j}(\theta)=\theta_j$ for $j\in J$. All remaining off-diagonal entries (\emph{completion}
entries) have cost $\Gamma$. The scenario box is $\Theta=\prod_j[0,\beta_j]$ and $F$ is the identity.

The parameter $\theta_j$ occurs exactly in the entries $(f_j,f_j')$ and $(g_j,g_j)$, with coefficient $1$; the
second lies on $F$. The vertices $f_j,f'_j,g_j$ $(j\in J)$ are pairwise distinct, which gives (C2) and (C3); (C4) is
clear. The pattern of the matrix depends only on $D$ and $J$, not on the numbers $\alpha,\beta$.

\subsection{Proof of Theorem \ref{thm:adapter}}

\begin{lemma}\label{lem:paths}
$D'$ is acyclic, and every new vertex has exactly one entering and one leaving arc in $D'$. Contracting the chains
gives a bijection $P'\mapsto P$ between the $s$--$t$ paths of $D'$ and those of $D$; the path $P'$ passes through
$f_j,f_j'$ if and only if $a_j\in P$, and through $g_j,g_j'$ if and only if $b_j\in P$.
\end{lemma}
\begin{proof}
$D'$ is a subdivision of $D$, and the new vertices are the interior vertices of the chains.
\end{proof}

\begin{lemma}\label{lem:perm}
Let $\sigma\ne\mathrm{id}$ be a permutation of the vertices such that $(v,\sigma(v))$ is a designated entry whenever
$\sigma(v)\ne v$. Then the non-fixed points of $\sigma$ form one cycle, consisting of $e$ and an $s$--$t$ path of
$D'$. Conversely every $s$--$t$ path $P'$ of $D'$ defines such a permutation; we denote it by $\sigma_P$, where $P$ is
the path of $D$ corresponding to $P'$ by Lemma~\ref{lem:paths}.
\end{lemma}
\begin{proof}
Each nontrivial cycle of $\sigma$ is a directed cycle of $D'+e$. As $D'$ is acyclic it contains $e$; as the cycles
of $\sigma$ are disjoint there is only one.
\end{proof}

\begin{lemma}[cost identity]\label{lem:cost}
For every $s$--$t$ path $P$ of $D$ and every $\theta\in\Theta$,
\[
 C(\theta)\cdot(\sigma_P-F)=\varrho+\sum_{a\in P}\alpha_a
   +\sum_{j\in J}\bigl(\theta_j[a_j\in P]+(\beta_j-\theta_j)[b_j\in P]\bigr)\ \in\ [\varrho,\ \varrho+A_0+B_0].
\]
\end{lemma}
\begin{proof}
The off-diagonal entries of $\sigma_P$ are $e$ (cost $\varrho$) and the arcs of the chains of the arcs $a\in P$: the
first arc of each chain (cost $\alpha_a$), arcs of cost $0$, the arcs $f_j\to f_j'$ with $a_j\in P$ (cost
$\theta_j$) and the arcs $g_j\to g_j'$ with $b_j\in P$ (cost $\beta_j$), by Lemma~\ref{lem:paths}. Passing from $F$
to $\sigma_P$ removes the diagonal entries of the vertices on the cycle; these are $0$ except $\theta_j$ at $g_j$ for
$b_j\in P$. For the bounds note $0\le\theta_ju+(\beta_j-\theta_j)v\le\beta_j$ for $u,v\in\{0,1\}$.
The equality is algebraic: it holds for every real $\theta$ and without any assumption on signs; the box and
$\beta\ge0$ are used only for the bounds and, later, for the minimum.
\end{proof}

\begin{lemma}[completion]\label{lem:compl}
If a permutation $\sigma$ uses a completion entry, then $C(\theta)\cdot(\sigma-F)\ge\Gamma-B_0\ge\varrho+A_0+B_0+1$
for every $\theta\in\Theta$.
\end{lemma}
\begin{proof}
All entries are nonnegative, so $C(\theta)\cdot\sigma\ge\Gamma$, while $C(\theta)\cdot F=\sum_j\theta_j\le B_0$.
\end{proof}

\begin{lemma}[dual certificate]\label{lem:dual}
For $\theta\in\Theta$ put $\lambda_{g_j}=\theta_j$ for $j\in J$, $\lambda_v=0$ for all other rows $v$, and $\mu_v=0$ for
all columns. Then $\lambda_u+\mu_v\le C_{uv}(\theta)$ for all $u,v$, with equality on the diagonal, and
$\sum_v\lambda_v+\sum_v\mu_v=C(\theta)\cdot F$. Hence $F$ is optimal in every scenario.
\end{lemma}
\begin{proof}
In row $g_j$ the entries are $\theta_j$ (diagonal), $\beta_j\ge\theta_j$ (the arc $g_j\to g'_j$) and
$\Gamma>\beta_j$ (completion). In all other rows the left side is $0$ and the entries are nonnegative.
\end{proof}

\begin{proof}[Proof of Theorem \ref{thm:adapter}]
By Lemmas \ref{lem:perm}--\ref{lem:compl} every permutation $\sigma\ne\mathrm{id}$ satisfies
$C(\theta)\cdot(\sigma-F)\ge\varrho>0$. By Birkhoff's theorem every $X\in\cA_N$ is a convex combination of
permutation matrices; if $X\ne F$ one of them is not the identity, so $C(\theta)\cdot(X-F)>0$. This proves
uniqueness in (C1), and Lemma~\ref{lem:dual} gives the dual certificate.

Let $X\in\cA_N$ with $X_e=1$. Every permutation in a Birkhoff decomposition of $X$ has $e$-entry $1$. Such a
permutation either uses a completion entry, and then costs at least $\varrho+A_0+B_0+1$ more than $F$
(Lemma~\ref{lem:compl}), or uses designated entries only, and then equals some $\sigma_P$ (Lemma~\ref{lem:perm}),
whose relative cost is given by Lemma~\ref{lem:cost} and is at most $\varrho+A_0+B_0$. Since $D$ has an $s$--$t$
path, the minimum over the face $X_e=1$ is attained by some $\sigma_P$, which is the formula for $h_e(\theta)$.

Finally $\min_\theta\min_P=\min_P\min_\theta$, the parameters are independent, and
$\min_{0\le\theta\le\beta}\{\theta u+(\beta-\theta)v\}$ equals $0$ unless $u=v=1$, when it equals $\beta$. This gives
$\eta_e=\varrho+\min_PQ(P)$.
\end{proof}

\begin{remark}\label{rem:gap}
The proof shows $C(\theta)\cdot(X-F)\ge\varrho$ on the face $X_e=1$ and for every permutation $X\ne F$. There is no
uniform positive gap over fractional points close to $F$. Outside the box $\Theta$ the identity need not be optimal.
The formalisation of Section~\ref{sec:formal} records which hypothesis each step uses: Lemma~\ref{lem:dual} needs
neither a path nor acyclicity nor $\varrho>0$, only nonnegativity, the box and $B_0\le\Gamma$; $\varrho>0$ is used
only for strict uniqueness; the reduction of the face $X_e=1$ to permutations through $e$ holds for every cost
matrix; and the assumption $a_j\ne b_j$ is not used at all.
\end{remark}

\begin{proof}[Proof of Corollary \ref{cor:inapx}]
Deciding whether a directed acyclic graph $G$ has an $s$--$t$ path that contains at most one vertex from each of
a given list of pairs of distinct vertices is NP-complete \cite{GMO}; see also \cite{Kovac}. We may assume that $G$ has an
$s$--$t$ path. Replace every vertex $v$ of $G$ by two vertices $v^-,v^+$ and an arc $\bar a_v=(v^-,v^+)$, and
every arc $(u,v)$ of $G$ by $(u^+,v^-)$. The paths from $s^-$ to $t^+$ in the new directed acyclic graph correspond
to the $s$--$t$ paths of $G$, and such a path contains $v$ if and only if it uses $\bar a_v$. Let $\alpha\equiv0$
and $\varrho=1$, and for every listed pair $\{u,v\}$ introduce one interaction $(\bar a_u,\bar a_v)$ with
coefficient $M$, a positive integer. By Theorem~\ref{thm:adapter}, $\eta_e=1$ if a path of the required kind
exists and $\eta_e\ge1+M$ otherwise. The dimension $N$ does not depend on $M$; take $M=2^{p(N)}$ and
$\Gamma=2+2M|J|$, so that all entries are integers with polynomially many digits. An algorithm as in the statement
would return a number at most $M$ in the first case and at least $1+M$ in the second.
\end{proof}

\section{The diamond family and the correlation polytope}\label{sec:diamond}

Let $D_n$ have vertices $v_0,\dots,v_n$ and $b_{i,0},b_{i,1}$ $(1\le i\le n)$ and arcs
$d_{i,a}=(v_{i-1},b_{i,a})$, $(b_{i,a},v_i)$; $s=v_0$, $t=v_n$. The $s$--$t$ paths are the paths $P_z$, $z\in\{0,1\}^n$, where $P_z$ uses $d_{i,z_i}$ in layer
$i$. Give
$d_{i,a}$ the cost $w_{i,a}$, the other arcs cost $0$, and include the interaction $(d_{i,a},d_{j,b})$ with
coefficient $w_{ij,ab}$ for all $i<j$, $a,b\in\{0,1\}$, whether or not the weight vanishes. There are $4n$ arcs,
$3n+1$ vertices and $2n(n-1)$ interactions, so $N=(3n+1)+4n+8n(n-1)=8n^2-n+1$, and $Q(P_z)=w\cdot m_n(z)$.
Theorem~\ref{thm:adapter} with $\varrho=1$ gives \eqref{eq:eta}. The map from $w$ to the pair (matrix of lower
endpoints, matrix of upper endpoints) is affine.

\begin{remark}\label{rem:const}
Condition \eqref{eq:Gamma} allows $\Gamma(w)=2+\sum w_{i,a}+2\sum w_{ij,ab}$, affine in $w$, or, for weights in
$[0,5]$, the constant $\Gamma=2+10n+20n(n-1)$.
\end{remark}

Each $m_n(z)$ has exactly one coordinate equal to $1$ in each of the $n$ singleton blocks $\{(i,0),(i,1)\}$ and in
each of the $\binom n2$ pair blocks $\{(ij,ab)\}_{a,b}$; so $\one\cdot m_n(z)=\kappa_n:=n+\binom n2$. Let
$Q_n=\operatorname{conv}\{m_n(z)\}$. In the coordinates $x_i=m_{i,1}$, $X_{ij}=m_{ij,11}$ one has
$m_{i,0}=1-x_i$, $m_{ij,10}=x_i-X_{ij}$, $m_{ij,01}=x_j-X_{ij}$, $m_{ij,00}=1-x_i-x_j+X_{ij}$; so $Q_n$ is affinely
isomorphic to the correlation polytope $\COR(n)=\operatorname{conv}\{(z_i,z_iz_j)\}$, and $\xc(Q_n)=\xc(\COR(n))$.

\begin{proof}[Proof of Corollary \ref{cor:np}]
Let $G$ be a graph on $[n]$. Put all weights equal to $1$ except $w_{ij,00}=w_{ij,11}=5$ for the edges $ij$ of $G$.
Then $\eta_n(w)=1+\kappa_n+4\,u(G)$, where $u(G)$ is the minimum number of edges of $G$ that are not cut by a
bipartition of $[n]$, and the instance has polynomial size. Maximum cut is NP-hard already for unweighted graphs \cite{GJS}; the
weighted problem is in Karp's list \cite{Karp}.
\end{proof}

\section{Exact representations}\label{sec:exact}

We use two standard facts. (i) If a linear function is bounded below on a nonempty polyhedron, its minimum is
attained. (ii) (\emph{Affine Farkas lemma.}) If $P_0=\{y:Gy\ge g,\ G'y=g'\}$ is nonempty and the affine function
$\ell$ is nonnegative on $P_0$, then $\ell=\lambda_0+\sum_i\lambda_i(G_iy-g_i)$ on $P_0$ for some $\lambda\ge0$; if
moreover $\ell$ vanishes at a point of $P_0$, then $\lambda_0=0$. A face of a polyhedron is described by the
inequalities of the polyhedron with some of them turned into equations, so it needs no more inequalities.

\begin{lemma}[values force a face]\label{lem:face}
Let $Q\subseteq\R^d$ be the convex hull of a nonempty set of $0/1$ vectors with the same coordinate sum $\kappa$. Let
$(P,a,\pi)$ satisfy $\min_{y\in P}\{a(y)+w\cdot\pi(y)\}=1+\min_{q\in Q}w\cdot q$ for all $w\in\Z^d_{\ge0}$. Then
$P_0=\{y\in P:a(y)+\one\cdot\pi(y)=1+\kappa\}$ is a nonempty face of $P$ with $a=1$ on $P_0$ and $\pi(P_0)=Q$. In
particular $\xc(Q)\le\size(P)$.
\end{lemma}
\begin{proof}
The weight $w=0$ gives $a\ge1$ on $P$. Fix $y\in P$ and $w\in\Z^d_{\ge0}$. For every positive integer $\lambda$,
$a(y)+\lambda w\cdot\pi(y)\ge1+\lambda\min_Qw\cdot q$; dividing by $\lambda$ and letting $\lambda\to\infty$ gives
$w\cdot\pi(y)\ge\min_Qw\cdot q$. By scaling this holds for rational, and by continuity for all real $w\ge0$. The set
$Q^\uparrow:=Q+\R^d_{\ge0}$ is closed, convex and upward closed, hence the intersection of the half-spaces
$\{x:w\cdot x\ge\min_Qw\cdot q\}$, $w\ge0$. So $\pi(P)\subseteq Q^\uparrow$, and in particular $\one\cdot\pi\ge\kappa$
on $P$.

Let $q$ be a vertex of $Q$ and $w_q=2\cdot\one-q$. For $p\in Q$ we have $w_q\cdot p=2\kappa-q\cdot p\ge\kappa$, with
equality only for $p=q$ (as $0\le p\le\one$ and $\one\cdot p=\kappa$); since $w_q>0$, $q$ is the unique minimiser of
$w_q$ over $Q^\uparrow$. Let $y_q$ minimise $a+w_q\cdot\pi$ over $P$; the value is $1+\kappa$, while $a(y_q)\ge1$ and
$w_q\cdot\pi(y_q)\ge\kappa$. Hence $a(y_q)=1$ and $\pi(y_q)=q$.

The weight $w=\one$ shows that $a+\one\cdot\pi\ge1+\kappa$ is valid on $P$, so $P_0$ is a face; it contains every
$y_q$. On $P_0$ both $a\ge1$ and $\one\cdot\pi\ge\kappa$ are tight, so $a=1$, and $\pi(y)\in Q^\uparrow$ with
coordinate sum $\kappa$ lies in $Q$. Thus $\pi(P_0)\subseteq Q$ is convex and contains all vertices of $Q$.
\end{proof}

\begin{proof}[Proof of Theorem \ref{thm:exact}]
By \eqref{eq:eta} and Lemma~\ref{lem:face} with $Q=Q_n$, every representation has size at least
$\xc(Q_n)=\xc(\COR(n))$. Conversely, if $Q_n=\pi(P)$ for a polyhedron $P$, then $(P,1,\pi)$ is a representation.
The bounds $1.5^n\le\xc(\COR(n))\le2^n$ are \cite{KW} and the simplex of distributions on $\{0,1\}^n$.
\end{proof}

\begin{proof}[Proof of Corollary \ref{cor:tw}]
The instance is that of Section~\ref{sec:diamond} with the interactions $(d_{i,a},d_{j,b})$ kept only for
$ij\in E(G)$. It has $3n+1$ vertices, $4n$ arcs and $4|E(G)|$ interactions, hence $N_G=7n+1+16|E(G)|$, and
Theorem~\ref{thm:adapter} gives \textup{(C1)--(C4)} and $\eta_G(w)=1+\min_zw\cdot m_G(z)$, where $m_G(z)$ consists
of the singleton blocks of $m_n(z)$ and of its pair blocks for $ij\in E(G)$. The vectors $m_G(z)$ are $0/1$ vectors
with coordinate sum $n+|E(G)|$, and $Q_G=\operatorname{conv}\{m_G(z)\}$ is affinely isomorphic to $\COR(G)$ by the
change of coordinates of Section~\ref{sec:diamond}. Lemma~\ref{lem:face} with $Q=Q_G$ and the converse argument in
the proof of Theorem~\ref{thm:exact} give the first statement. The bounds in terms of the treewidth are Theorems~1
and~3 of \cite{AFHMS}. For the last statement, the stable set polytope of $G$ is, after deleting these zero coordinates, the face of
$\COR(G)$ on which all coordinates $z_iz_j$, $ij\in E(G)$, vanish, so $\xc(\COR(G))$ is at least the extension complexity of the stable set
polytope; by \cite[Theorem~1]{GJW} this is $2^{\Omega(n/\log n)}$ for an explicit family of $n$-vertex graphs of
bounded degree, for which $|E(G)|=O(n)$ and hence $N_G=\Theta(n)$.
\end{proof}

\subsection*{Bounded weights}
A set $R$ of pairs $(S,T)$ of disjoint subsets of $[n]$ is \emph{valid} if $|S\cap T'|\ne1$ for all
$(S,T),(S',T')\in R$. Kaibel and Weltge prove \cite[proof of Theorem~1]{KW}: \emph{the set of all pairs of disjoint
subsets of $[n]$ cannot be covered by fewer than $1.5^n$ valid sets.}

\begin{proof}[Proof of Theorem \ref{thm:bounded}]
Write $h(w)=\min_{Q_n}w\cdot q$ and $\kappa=\kappa_n$. Since $\one\cdot q=\kappa$ on $Q_n$, we have
$\eta_n(w+\one)=1+\kappa+h(w)$. So if $(P,a,\pi)$ represents $\eta_n$ on $\{1,\dots,5\}^d$, then
$(P,\,a+\one\cdot\pi-\kappa,\,\pi)$ has value $1+h(w)$ for all $w\in\{0,\dots,4\}^d$. Renaming, assume
\begin{equation}\label{eq:rep04}
 \min_{y\in P}\{a(y)+w\cdot\pi(y)\}=1+h(w)\qquad(w\in\{0,1,2,3,4\}^d).
\end{equation}

\emph{Step 1: exact lifts.} From $w=0$ and $w=\one$: $a\ge1$ and $a+\one\cdot\pi\ge1+\kappa$ on $P$. Let $q$ be a
vertex of $Q_n$, $w_q=2\cdot\one-q$, and let $y_q$ minimise $a+w_q\cdot\pi$, so $a(y_q)+w_q\cdot\pi(y_q)=1+\kappa$.
Put $\tau=a(y_q)-1\ge0$. The weight $2w_q$ is allowed and $h(2w_q)=2\kappa$, so
$1+\tau+2(\kappa-\tau)\ge1+2\kappa$, that is $\tau=0$. A unit vector as weight gives $a+\pi_c\ge1$ for every
coordinate $c$, and the indicator of a block $B$ gives $a+\sum_{c\in B}\pi_c\ge2$. At $y_q$, where $a=1$: $\pi\ge0$
and every block sum is at least $1$. As $w_q$ equals $1$ on the support of $q$ and $2$ elsewhere,
$\kappa=w_q\cdot\pi(y_q)=\sum_B\sum_{c\in B}\pi_c(y_q)+\sum_{c\notin\operatorname{supp}q}\pi_c(y_q)\ge\kappa$, which
forces $\pi(y_q)=q$.

\emph{Step 2: a face.} $P_0=\{y\in P:2a(y)+\one\cdot\pi(y)=2+\kappa\}$ is a face of $P$ containing all $y_q$. On
$P_0$: $a=1$, $\one\cdot\pi=\kappa$, $\pi\ge0$, and every block sum equals $1$.

\emph{Step 3: a slack matrix.} For $\emptyset\ne S\subseteq[n]$ let $w_S$ have coefficient $1$ at $(i,0)$ for
$i\in S$, $2$ at $(ij,11)$ for $i<j$ in $S$, and $0$ elsewhere. For $T\subseteq[n]$, $q_T=m_n(\one_T)$ and
$k=|S\cap T|$,
\[
 w_S\cdot q_T=(|S|-k)+k(k-1)=|S|-1+(k-1)^2,\qquad h(w_S)=|S|-1 .
\]
Hence $\ell_S:=w_S\cdot\pi-(|S|-1)$ is nonnegative on $P_0$, and $\ell_S(y_T)=(|S\cap T|-1)^2$ where $y_T=y_{q_T}$.
Let $s_1,\dots,s_r$ be the slacks of the inequalities of $P_0$; $r\le\size(P)$. Since $\ell_S$ vanishes at some
$y_T$, the affine Farkas lemma gives $\ell_S=\sum_i\lambda_{S,i}s_i$ on $P_0$ with $\lambda_{S,i}\ge0$. Likewise each
coordinate function $\pi_c\ge0$ vanishes at some $y_T$, so $1=\kappa^{-1}\sum_c\pi_c$ is a nonnegative combination
$\sum_i\lambda_{\emptyset,i}s_i$ on $P_0$. Therefore the matrix $M_{S,T}=(|S\cap T|-1)^2$, $S,T\subseteq[n]$, equals
$\sum_{i=1}^r\lambda_{S,i}\,s_i(y_T)$.

\emph{Step 4: covering.} Let $R_i$ be the set of disjoint pairs $(S,T)$ with $\lambda_{S,i}s_i(y_T)>0$. If
$(S,T),(S',T')\in R_i$ then $\lambda_{S,i}s_i(y_{T'})>0$, so $M_{S,T'}>0$ and $|S\cap T'|\ne1$: $R_i$ is valid. Every
disjoint pair has $M_{S,T}=1$ and lies in some $R_i$. By the result quoted above, $r\ge1.5^n$.
\end{proof}

\begin{remark}\label{rem:letters}
The proof uses only the weight values $0,1,2,4$, so the theorem holds for the alphabet $\{1,2,3,5\}$.
\end{remark}

\begin{remark}[weights in the right-hand side]\label{rem:rhs}
Let $A,A',B,B'$ be matrices and $b,b',c$ vectors, and suppose
$\eta(w)=\max\{c\cdot x:\ Ax\le b+Bw,\ A'x=b'+B'w\}$ for all $w\in W$, with $m$ inequality rows. By
linear programming duality $\eta(w)=\min\{(b+Bw)\cdot u+(b'+B'w)\cdot v:\ A^{\top}u+A'^{\top}v=c,\ u\ge0\}$, and
this is a weight-affine LP representation with $m$ inequalities. Hence Theorems
\ref{thm:exact}--\ref{thm:abs} hold with the same bounds for maximisation programs in which $w$ enters only the
right-hand side. Conversely, the dual of a weight-affine LP is a maximisation program in which $w$ occurs only in
the right-hand side and in an additive affine term. A minimisation program with $w$ only in the right-hand side has
a convex value function and therefore cannot represent $\eta_n$ on any set $W$ that contains, for one singleton
block, the weights $(1,3),(2,2),(3,1)$ with all other weights equal to $1$, in particular on the sets $W$ of
Theorems \ref{thm:exact}--\ref{thm:abs}: the values there are $v,v+1,v$. On a two-letter alphabet this obstruction
disappears.
\end{remark}

\begin{proof}[Proof of Proposition \ref{prop:upper}]
Let $e=(r,c)$. For fixed $\theta$, the minimum of $C(\theta)\cdot X$ over $X\in\cA_N$ with $X_e=1$ is $C_e(\theta)$ plus
the value of the assignment problem on the rows other than $r$ and the columns other than $c$, which by duality
equals $\max\{\sum_{i\ne r}u_i+\sum_{j\ne c}v_j:\ u_i+v_j\le C_{ij}(\theta)\text{ for }i\ne r,\ j\ne c\}$. Hence
\[
 \max_{\theta\in\Theta}h_e(\theta)=\max\Bigl\{C_e(\theta)-C(\theta)\cdot F+\sum_{i\ne r}u_i+\sum_{j\ne c}v_j:\
 u_i+v_j\le C_{ij}(\theta)\ (i\ne r,\ j\ne c),\ \theta\in\Theta\Bigr\},
\]
a linear program in $(\theta,u,v)$, because $C$ is affine. For the instances of Theorem~\ref{thm:adapter}, the
theorem says that $h_e(\theta)-\varrho$ is the length of a shortest $s$--$t$ path of $D$ for the arc lengths
$\ell_a(\theta)\ge0$, which is $\max\{\pi_t-\pi_s:\pi_y-\pi_x\le\ell_a(\theta)\text{ for every arc }a=(x,y)\}$;
maximising over $\theta$ as well gives the displayed program, which has $|A|+2|J|$ inequalities. In the family
\eqref{eq:eta} there are $4n$ arcs and $2n(n-1)$ interactions, and the weights occur only in the right-hand sides
($\alpha$ and $\beta$ are coordinates of $w$). By Remark~\ref{rem:rhs} the maximum has a weight-affine LP
representation with $4n+4n(n-1)=4n^2$ inequalities.
\end{proof}

\section{Fixed additive error}\label{sec:abs}

We use the following lemma \cite[Lemma~4]{BFPS}. Let $u\equiv3\pmod4$, $\ell=(u+1)/4$, and let $\mathcal D$ and
$\mathcal U$ be the sets of pairs of $\ell$-subsets of $[u]$ with intersection of size $0$ and $1$. \emph{For all
nonnegative functions $\phi,\psi$ on subsets of $[u]$, $X(S,T):=\phi(S)\psi(T)$, and $0<\xi<1$,}
\begin{equation}\label{eq:bfps}
 (1-\xi)\,\E[X\mid\mathcal D]-\E[X\mid\mathcal U]\ \le\ \|X\!\restriction\!(\mathcal D\cup\mathcal U)\|_\infty\;
   2^{-\frac{\xi^2}{16\ln2}\ell+O(\log\ell)},
\end{equation}
\emph{with an absolute constant in $O(\log\ell)$; expectations are uniform over $\mathcal D$ and $\mathcal U$.}

\begin{proposition}\label{prop:small}
Let $n\ge3$ and let $(P,a,\pi)$ have value function $g$ with $|g(w)-\eta_n(w)|\le\frac1{40}$ on
$\{1,\dots,5\}^{d_n}$. Then $\size(P)\ge
2^{\ell/(64\ln2)-O(\log(\ell+1))}$ with $\ell=\lfloor(n+1)/4\rfloor$.
\end{proposition}
\begin{proof}
As in the proof of Theorem~\ref{thm:bounded}, replacing $a$ by $b:=a+\one\cdot\pi-\kappa-\frac1{40}$ we get, with
$\gamma=\frac1{20}$ and $f(w):=\min_P\{b+w\cdot\pi\}$,
\begin{equation}\label{eq:sandwich}
 1+h(w)-\gamma\le f(w)\le1+h(w)\qquad(w\in\{0,\dots,4\}^d).
\end{equation}

\emph{Approximate lifts.} Let $q$ be a vertex of $Q_n$ and $c_q=3\cdot\one-2q$, so $h(c_q)=\kappa$. Let $y_q$ minimise
$b+c_q\cdot\pi$ and write $b(y_q)+c_q\cdot\pi(y_q)=1+\kappa-s_q$ with $0\le s_q\le\gamma$. For a sign vector
$\sigma\in\{\pm1\}^d$ the weight $c_q+\sigma$ is allowed; in each block the coordinate of $q$ costs at most $2$ and
the others at least $2$, so $h(c_q+\sigma)=\kappa+\sigma\cdot q$. The lower bound in \eqref{eq:sandwich} at
$c_q+\sigma$, evaluated at $y_q$, gives $\sigma\cdot\zeta_q\ge s_q-\gamma$ for $\zeta_q:=\pi(y_q)-q$. Choosing
$\sigma$ opposite to the signs of $\zeta_q$,
\begin{equation}\label{eq:lift}
 \|\pi(y_q)-q\|_1\le\gamma-s_q,\qquad b(y_q)-1=-s_q-c_q\cdot\zeta_q .
\end{equation}

\emph{A perturbed matrix.} With $w_S$ as before, $L_S:=b+w_S\cdot\pi-|S|+\gamma$ is nonnegative on $P$ by
\eqref{eq:sandwich}. Put $M'_{S,T}=L_S(y_T)$. By \eqref{eq:lift}, writing
$s_T,\zeta_T$ for $s_{q_T},\zeta_{q_T}$ and $v_T=\gamma-s_T\in[0,\gamma]$,
\[
 M'_{S,T}=(|S\cap T|-1)^2+v_T+(w_S-c_{q_T})\cdot\zeta_T,\qquad |(w_S-c_{q_T})\cdot\zeta_T|\le3v_T .
\]
So $1-2\gamma\le M'_{S,T}\le1+4\gamma$ if $S\cap T=\emptyset$ and $0\le M'_{S,T}\le4\gamma$ if $|S\cap T|=1$. By the
affine Farkas lemma each $L_S$ is a nonnegative combination of the $r=\size(P)$ slacks of $P$ and the constant $1$;
hence $M'$ is a sum of $r+1$ functions of the form $\phi(S)\psi(T)\ge0$, each pointwise at most $M'$.

\emph{Corruption.} Restrict to $\ell$-subsets of $[u]$, $u=4\ell-1\le n$, and apply \eqref{eq:bfps} with
$\xi=\frac12$ to each of the $r+1$ terms; on $\mathcal D\cup\mathcal U$ each is at most $1+4\gamma=\frac65$. Summing,
$\frac14=\frac12(1-2\gamma)-4\gamma\le\frac12\E[M'\mid\mathcal D]-\E[M'\mid\mathcal U]\le(r+1)\,\frac65\,2^{-\ell/(64\ln2)+O(\log\ell)}$.
Finally $r\ge1$: a program without inequalities has an affine value function, and $\eta_n$ is not within
$\frac1{40}$ of an affine function (let one singleton block run through $(1,3),(2,2),(3,1)$ with all other weights
equal to $1$; the values are $v,v+1,v$). So $r+1\le2r$, and the constants are absorbed in $O(\log(\ell+1))$.
\end{proof}

\begin{proof}[Proof of Theorem \ref{thm:abs}]
Split $\nu m$ of the $n$ variables into $\nu$ groups of size $m$. For $w\in\{1,\dots,5\}^{d_m}$ let $W(w)\in
\{1,\dots,5\}^{d_n}$ copy $w$ inside each group and equal $1$ on all other blocks; there are $K=\kappa_n-\nu\kappa_m$
such blocks, each contributing exactly $1$. Then $\eta_n(W(w))=\nu\,\eta_m(w)+K+1-\nu$. Since $W$ is affine,
$\hat g(w):=(g(W(w))-K-1+\nu)/\nu$ is the value of a weight-affine LP on the same polyhedron, with error at most
$\Delta/\nu\le\frac1{40}$. Apply Proposition~\ref{prop:small} with $m$ in place of $n$.
\end{proof}

\section{Fixed relative error}\label{sec:rel}

\begin{proof}[Proof of Theorem \ref{thm:rel}]
\emph{The linear program.} Let $r=t+1$. Variables are numbers $\mu_S(x_S)$ for $S\subseteq[n]$, $|S|\le r$,
$x_S\in\{0,1\}^S$, subject to the equations $\sum_{x_S}\mu_S(x_S)=1$ and $\mu_S=$ marginal of $\mu_{S'}$ for
$S\subseteq S'$, and to the inequalities $\mu_S(x_S)\ge0$ for $|S|=r$ (nonnegativity for smaller $S$ follows). There
are $2^r\binom nr$ inequalities. Let $L_t(w)=1+\min_\mu\bigl[\sum_{i,a}w_{i,a}\mu_i(a)+\sum_{i<j,a,b}w_{ij,ab}\mu_{ij}(a,b)\bigr]$.
Every $z$ gives a feasible point, so $L_t\le\eta_n$; and as each block of marginals is a probability vector and
$w\ge1$, $L_t\ge1+\kappa_n$.

\emph{Entropy.} Fix a feasible $\mu$. For a tuple $S$ of at most $t$ indices (repetitions allowed) and an index $i$,
the conditional entropy $\Ent(X_i\mid X_S)$ is computed in the local distribution on the set of indices involved; so
is the conditional mutual information $\MI(X_i;X_j\mid X_S)$ when $S$ has at most $t-1$ entries. In both cases at
most $t+1$ indices are involved, and consistency makes this well defined. For $0\le s\le t$ let $H_s=\E\,\Ent(X_i\mid X_{S})$, where $S$ is a uniform $s$-tuple and $i$ a uniform
index, independent. The chain rule gives $H_s-H_{s+1}=\E\,\MI(X_i;X_j\mid X_S)$ for $s<t$, and
$0\le H_t\le H_0\le\ln2$. Hence for some $s<t$,
\begin{equation}\label{eq:mi}
 \E_{S,i,j}\,\MI(X_i;X_j\mid X_S)\le\frac{\ln2}{t}.
\end{equation}

\emph{Rounding.} For this $s$ draw $S$, then $x_S$ from $\mu_S$, and then draw all $Z_i$ independently from the
conditional laws $p_i$ of $X_i$ given $X_S=x_S$. For $i\ne j$ let $p_{ij}$ be the conditional law of $(X_i,X_j)$. By
Pinsker's inequality $\TV(p_{ij},p_i\otimes p_j)\le\sqrt{\MI(X_i;X_j\mid X_S=x_S)/2}$. A pair block costs between
$1$ and $5$, so replacing $p_{ij}$ by $p_i\otimes p_j$ changes its expected cost by at most $4\TV$. Averaging over
the seed, summing over pairs and using the Cauchy--Schwarz inequality with \eqref{eq:mi},
\[
 \E\sum_{i<j}4\,\TV(p_{ij},p_i\otimes p_j)\le2n^2\sqrt{\frac{\ln2}{2t}}=n^2\sqrt{\frac{2\ln2}{t}} .
\]
Averaging the conditional laws $p_i$, $p_{ij}$ over the seed returns $\mu_i$, $\mu_{ij}$. So the expected cost of
$Z$ is at most the objective value of $\mu$ plus the last bound, and some $z$ is at least as good. Applying this to
an optimal $\mu$ proves the first display. Since $1+\kappa_n\ge n^2/2$, the additive error is at most
$\sqrt{8\ln2/t}\,L_t$, and $t=\lceil8\ln2/\varepsilon^2\rceil$ gives the second statement.
\end{proof}

\begin{remark}
The pair-local relaxation $(t=1)$ is not exact: with all weights $1$ except $5$ on equal pair states, its value is
$1+\kappa_n$, while $\eta_n=1+\kappa_n+n^2-2n$ for even $n\ge4$. The size bound of Theorem~\ref{thm:rel} is below $2^n$
only for fairly large $n$ (from $n=27$ for $\varepsilon=1$).
\end{remark}

\section{Formal verification and reproducibility}\label{sec:formal}
A Lean~4 project (Lean 4.19.0, Mathlib commit \code{c44e0c8e}, both pinned in the project; one check command; a
recorded clean-room run) accompanies the paper. It contains twelve modules, $3907$ lines. No \code{sorry} occurs, and
the $57$ reported theorems depend only on Lean's three standard axioms.

\emph{Theorem~\ref{thm:adapter} is machine-checked.} The theorem \code{exact\_adapter} assumes finite sets $V$,
$A$, $J$, maps giving the tail and head of every arc and the two arcs of every interaction, $s\ne t$, a rank
function increasing along every arc, at least one $s$--$t$ path, $\alpha\ge0$, $\beta\ge0$, $\varrho>0$ and
\eqref{eq:Gamma}. It concludes, for the cost matrix of Section~\ref{sec:adapter} defined on all ordered pairs of
vertices:
\begin{itemize}
\item the dimension is $|V|+|A|+4|J|$;
\item \textup{(C2)}: $C(\theta)=C(0)+\sum_j\theta_j\,(E_{f_jf_j'}+E_{g_jg_j})$ for every real $\theta$, where $E_{uv}$
  is the matrix unit; \textup{(C3)}: the $2|J|$ entries $(f_j,f_j')$, $(g_j,g_j)$ lie in pairwise different rows and
  pairwise different columns;
\item for every $\theta\in\Theta$: \textup{(C4)}; the dual certificate of Lemma~\ref{lem:dual}; and
  $C(\theta)\cdot X>C(\theta)\cdot F$ for every doubly stochastic $X\ne F$;
\item for every $\theta\in\Theta$: the set of values $C(\theta)\cdot(X-F)$ over the doubly stochastic $X$ with $X_e=1$
  and the set of numbers $\varrho+\sum_{a\in P}\alpha_a+\sum_j(\theta_j[a_j\in P]+(\beta_j-\theta_j)[b_j\in P])$ over the
  $s$--$t$ paths $P$ of $D$ have least elements, and these are equal;
\item the set of values $C(\theta)\cdot(X-F)$ over all $\theta\in\Theta$ and all such $X$, and the set of numbers
  $\varrho+Q(P)$, have least elements, and these are equal.
\end{itemize}
Paths of $D$ are lists of arcs, so parallel arcs are allowed. The proof follows Section~\ref{sec:adapter}:
permutations supported on designated entries (Lemma~\ref{lem:perm}), forced traversal of whole chains, the cost
identity of Lemma~\ref{lem:cost} for every real $\theta$, Lemmas \ref{lem:compl} and \ref{lem:dual}, and the
passage from permutations to doubly stochastic matrices by the Birkhoff--von Neumann theorem, which is taken from
Mathlib. Separate theorems state that the correspondence $P\mapsto\sigma_P$ of Lemmas \ref{lem:paths} and
\ref{lem:perm} is a bijection (for every $s$--$t$ path of $D$ there is exactly one permutation supported on
designated entries, through $e$, that traverses exactly the chains of its arcs, and conversely), and that the dual
prices are $\lambda_x=\sum_j\theta_j[x=g_j]$, hence linear in $\theta$. Two qualifications. First,
Lemma~\ref{lem:paths} is not formalised as written: Lean has no statement on $s$--$t$ paths of $D'$. It states the
bijection between the paths of $D$ and the permutations directly, and that such a permutation enters a chain if
and only if it moves all its new vertices. For new vertices, `at most one' entering and leaving arc is a Lean
statement; `exactly one' is not. Secondly, the agreement of the Lean definitions with the construction of
Section~\ref{sec:adapter} was checked by reading them against the text and has not been reviewed externally. What
is trusted is Lean's kernel, the three axioms, and Mathlib as distributed: the kernel replay covers the twelve
modules, not Mathlib.

\emph{Also formalised}, in the module \code{CorrelationCoverKernel}: the covering inequality of Kaibel and Weltge
\cite{KW} in the language of words over a three-letter alphabet (a valid family has at most $2^n$ members, hence
$r$ valid families covering all $3^n$ words satisfy $3^n\le r\,2^n$). This is a formalisation of a known argument.

\emph{Not formalised:} the diamond family of Section~\ref{sec:diamond}; Propositions \ref{prop:indep} and
\ref{prop:upper}; Corollaries \ref{cor:np}, \ref{cor:inapx} and \ref{cor:tw}; Lemma~\ref{lem:face}, the passage
from a linear program to a covering, and Theorems \ref{thm:exact}--\ref{thm:rel}.

\begin{sloppypar}
The Lean project and the scripts of the computations described in the introduction are in the package
accompanying this paper: the directory \code{forced-entry-costs} of the repository
\url{https://github.com/kosza123/flow-certificate-thresholds}, and the Zenodo record of this paper.
\end{sloppypar}

\section{Open questions}
\textup{(1)} The trade-off between additive error and size between the regimes of Theorems \ref{thm:abs} and
\ref{thm:rel}. \textup{(2)}
Whether the bound $2^{\Omega(N/\log N)}$ of Corollary~\ref{cor:tw} can be improved to $2^{\Omega(N)}$; this would
follow from \cite[Conjecture~2]{AFHMS} or from the conjecture of \cite{GJW}. \textup{(3)} Linear
programs in which the weights enter the constraint matrix, or the objective and the right-hand side together; the
right-hand side alone is covered by Remark~\ref{rem:rhs}. \textup{(4)} Theorem~\ref{thm:bounded} for an alphabet of
three letters.

\section*{Author contributions and use of AI}
This work was carried out with substantial assistance from AI systems working under the author's direction. The
SAPIO MORPHEUS project uses models from OpenAI and Anthropic.
\begin{itemize}
\item The construction, the theorems on linear programs, their first complete proofs and three of the Lean modules
  were produced by an AI operator of the project based on OpenAI models, with its own checking agents.
\item The independent reviews of the proofs and of this text with separate code, the literature comparison,
  Propositions \ref{prop:indep} and \ref{prop:upper}, Corollaries \ref{cor:np}, \ref{cor:inapx} and
  \ref{cor:tw}, Remark~\ref{rem:rhs}, the restructuring and this manuscript, the Lean formalisation of
  Theorem~\ref{thm:adapter} built on those three modules, and the Lean project with its clean-room record
  were produced by AI agents based on Anthropic's Claude.
\item The author selected the problem and the direction, set the standards of rigour, made the decisions, and takes
  responsibility for the content.
\end{itemize}
In line with common publication policies, AI systems are not listed as authors.

\begin{thebibliography}{99}
\bibitem{AFHMS} P.~Aboulker, S.~Fiorini, T.~Huynh, M.~Macchia, J.~Seif, \emph{Extension complexity of the correlation
  polytope}, Oper.\ Res.\ Lett.\ \textbf{47} (2019), 47--51; arXiv:1806.00541v2.
\bibitem{ABV} H.~Aissi, C.~Bazgan, D.~Vanderpooten, \emph{Complexity of the min-max and min-max regret assignment
  problems}, Oper.\ Res.\ Lett.\ \textbf{33} (2005), 634--640.
\bibitem{ABTW} D.~Avis, D.~Bremner, H.~R.~Tiwary, O.~Watanabe, \emph{Polynomial size linear programs for problems in
  P}, Discrete Appl.\ Math.\ \textbf{265} (2019), 22--39; arXiv:1408.0807.
\bibitem{BHKL} A.~E.~Black, C.~Hertrich, S.~Kober, G.~Loho, \emph{Two linear programs are better than one},
  arXiv:2610.15909 (2026).
\bibitem{BFPS} G.~Braun, S.~Fiorini, S.~Pokutta, D.~Steurer, \emph{Approximation limits of linear programs (beyond
  hierarchies)}, Math.\ Oper.\ Res.\ \textbf{40} (2015), 756--772; arXiv:1204.0957v3.
\bibitem{BP13} G.~Braun, S.~Pokutta, \emph{Common information and unique disjointness}, Algorithmica \textbf{76}
  (2016), 597--629.
\bibitem{BPmatching} G.~Braun, S.~Pokutta, \emph{The matching polytope does not admit fully-polynomial size
  relaxation schemes}, in: Proc.\ SODA 2015, 837--846; arXiv:1403.6710.
\bibitem{BPZ} G.~Braun, S.~Pokutta, D.~Zink, \emph{Affine reductions for LPs and SDPs}, Math.\ Program.\
  \textbf{173} (2019), 281--312; arXiv:1410.8816.
\bibitem{BM} M.~Braverman, A.~Moitra, \emph{An information complexity approach to extended formulations}, in:
  Proc.\ STOC 2013, 161--170.
\bibitem{Cela} E.~\c{C}ela, B.~Klinz, S.~Lendl, G.~J.~Woeginger, L.~Wulf, \emph{A linear time algorithm for
  linearizing quadratic and higher-order shortest path problems}, Math.\ Program.\ \textbf{210} (2025), 165--188.
\bibitem{CW} N.~Chakravarti, A.~P.~M.~Wagelmans, \emph{Calculation of stability radii for combinatorial optimization
  problems}, Oper.\ Res.\ Lett.\ \textbf{23} (1998), 1--7.
\bibitem{CLRS} S.~O.~Chan, J.~R.~Lee, P.~Raghavendra, D.~Steurer, \emph{Approximate constraint satisfaction requires
  large LP relaxations}, J.\ ACM \textbf{63} (2016), no.~4; arXiv:1309.0563.
\bibitem{DPSW} A.~Darmann, U.~Pferschy, J.~Schauer, G.~J.~Woeginger, \emph{Paths, trees and matchings under
  disjunctive constraints}, Discrete Appl.\ Math.\ \textbf{159} (2011), 1726--1735.
\bibitem{FMPTW} S.~Fiorini, S.~Massar, S.~Pokutta, H.~R.~Tiwary, R.~de~Wolf, \emph{Exponential lower bounds for
  polytopes in combinatorial optimization}, J.\ ACM \textbf{62} (2015), no.~2, Art.~17; arXiv:1111.0837.
\bibitem{Fortin} J.~Fortin, P.~Zieli\'nski, D.~Dubois, H.~Fargier, \emph{Criticality analysis of activity networks
  under interval uncertainty}, J.\ Sched.\ \textbf{13} (2010), 609--627.
\bibitem{GMO} H.~N.~Gabow, S.~N.~Maheshwari, L.~J.~Osterweil, \emph{On two problems in the generation of program
  test paths}, IEEE Trans.\ Software Eng.\ \textbf{SE-2} (1976), 227--231.
\bibitem{GJS} M.~R.~Garey, D.~S.~Johnson, L.~Stockmeyer, \emph{Some simplified NP-complete graph problems},
  Theoret.\ Comput.\ Sci.\ \textbf{1} (1976), 237--267.
\bibitem{GW} F.~Glover, E.~Woolsey, \emph{Converting the 0-1 polynomial programming problem to a 0-1 linear
  program}, Oper.\ Res.\ \textbf{22} (1974), 180--182.
\bibitem{GJM} B.~Goldengorin, G.~J\"ager, P.~Molitor, \emph{Tolerances applied in combinatorial optimization},
  J.\ Comput.\ Sci.\ \textbf{2} (2006), 716--734.
\bibitem{GJW} M.~G\"o\"os, R.~Jain, T.~Watson, \emph{Extension complexity of independent set polytopes}, SIAM J.\
  Comput.\ \textbf{47} (2018), 241--269; arXiv:1604.07062.
\bibitem{HL} C.~Hertrich, G.~Loho, \emph{Neural networks and (virtual) extended formulations}, SIAM J.\ Math.\ Data
  Sci.\ \textbf{8} (2026), 1350--1369; arXiv:2411.03006.
\bibitem{Hladik11} M.~Hlad\'{\i}k, \emph{Tolerance analysis in linear systems and linear programming}, Optim.\
  Methods Softw.\ \textbf{26} (2011), 381--396.
\bibitem{Hladik12} M.~Hlad\'{\i}k, \emph{Complexity of necessary efficiency in interval linear programming and
  multiobjective linear programming}, Optim.\ Lett.\ \textbf{6} (2012), 893--899.
\bibitem{Hladik26} M.~Hlad\'{\i}k, \emph{Tolerance analysis in linear programming: searching for lower and upper
  bounds}, European J.\ Oper.\ Res.\ \textbf{336} (2027), 39--52 (online 2026), doi:10.1016/j.ejor.2026.07.048.
\bibitem{HM} A.~J.~Hoffman, H.~M.~Markowitz, \emph{A note on shortest path, assignment, and transportation
  problems}, Naval Res.\ Logist.\ Quart.\ \textbf{10} (1963), 375--379.
\bibitem{HuSotirov} H.~Hu, R.~Sotirov, \emph{Special cases of the quadratic shortest path problem}, J.\ Comb.\
  Optim.\ \textbf{35} (2018), 754--777; arXiv:1611.07682.
\bibitem{JKR} V.~Jain, F.~Koehler, A.~Risteski, \emph{Mean-field approximation, convex hierarchies, and the
  optimality of correlation rounding: a unified perspective}, in: Proc.\ STOC 2019, 1226--1236; arXiv:1808.07226.
\bibitem{KW} V.~Kaibel, S.~Weltge, \emph{A short proof that the extension complexity of the correlation polytope
  grows exponentially}, Discrete Comput.\ Geom.\ \textbf{53} (2015), 397--401; corrected version arXiv:1307.3543v6.
\bibitem{Karp} R.~M.~Karp, \emph{Reducibility among combinatorial problems}, in: R.~E.~Miller, J.~W.~Thatcher (eds.),
  Complexity of Computer Computations, Plenum, 1972, 85--103.
\bibitem{KZ} A.~Kasperski, P.~Zieli\'nski, \emph{Minmax regret approach and optimality evaluation in combinatorial
  optimization problems with interval and fuzzy weights}, European J.\ Oper.\ Res.\ \textbf{200} (2010), 680--687.
\bibitem{Kovac} J.~Kov\'a\v{c}, \emph{Complexity of the path avoiding forbidden pairs problem revisited},
  Discrete Appl.\ Math.\ \textbf{161} (2013), 1506--1512; arXiv:1111.3996.
\bibitem{LinWen} C.-J.~Lin, U.-P.~Wen, \emph{Sensitivity analysis of the optimal assignment}, European J.\ Oper.\
  Res.\ \textbf{149} (2003), 35--46.
\bibitem{MR} P.~Manurangsi, P.~Raghavendra, \emph{A birthday repetition theorem and complexity of approximating
  dense CSPs}, in: Proc.\ ICALP 2017, LIPIcs \textbf{80}, Art.~78; arXiv:1607.02986.
\bibitem{MS} R.~Montemanni, D.~H.~Smith, \emph{On solving the shortest paths with exclusive-disjunction arc pairs
  conflicts}, arXiv:2506.03326 (2025).
\bibitem{RT} P.~Raghavendra, N.~Tan, \emph{Approximating CSPs with global cardinality constraints using SDP
  hierarchies}, in: Proc.\ SODA 2012, 373--387; arXiv:1110.1064.
\bibitem{Rostami} B.~Rostami, A.~Chassein, M.~Hopf, D.~Frey, C.~Buchheim, F.~Malucelli, M.~Goerigk, \emph{The
  quadratic shortest path problem: complexity, approximability, and solution methods}, European J.\ Oper.\ Res.\
  \textbf{268} (2018), 473--485.
\bibitem{SLG} Yu.~N.~Sotskov, V.~K.~Leontev, E.~N.~Gordeev, \emph{Some concepts of stability analysis in
  combinatorial optimization}, Discrete Appl.\ Math.\ \textbf{58} (1995), 169--190.
\bibitem{Wendell} R.~E.~Wendell, \emph{The tolerance approach to sensitivity analysis in linear programming},
  Management Sci.\ \textbf{31} (1985), 564--578.
\bibitem{Yannakakis} M.~Yannakakis, \emph{Expressing combinatorial optimization problems by linear programs},
  J.\ Comput.\ System Sci.\ \textbf{43} (1991), 441--466.
\bibitem{YZ} Y.~Yoshida, Y.~Zhou, \emph{Approximation schemes via Sherali--Adams hierarchy for dense constraint
  satisfaction problems and assignment problems}, in: Proc.\ ITCS 2014, 423--438.
\end{thebibliography}

\end{document}
